\chapter{Chern--Simons Theory}
\label{ch:tqft-chern-simons-theory}

\ChapterConventions{oriented smooth three-manifold \(M\)}
{no metric is part of the classical action}
{\(1\), unless explicitly displayed}
{a vacuum vector appears only after choosing a Hamiltonian splitting or a
boundary state space}
{gauge transformations act on connections by pullback and conjugation}

Let \(G\) be a compact Lie group with Lie algebra \(\mathfrak g\).  Fix an
invariant symmetric bilinear form \(\operatorname{tr}\) on \(\mathfrak g\).
For a connection \(A\) on a principal \(G\)-bundle over an oriented
three-manifold \(M\), the Chern--Simons functional in a local trivialization is
\[
  S_{\rm CS}(A)
  =
  \frac{k}{4\pi}\int_M
  \operatorname{tr}\!\left(A\wedge dA+\frac23 A\wedge A\wedge A\right).
\]
The level \(k\) is not a free real parameter once large gauge transformations
are included.

\section{Variation and Classical Equations}

Let \(F_A=dA+A\wedge A\).  For a compactly supported variation \(\delta A\)
on a closed manifold,
\[
  \delta S_{\rm CS}
  =
  \frac{k}{2\pi}\int_M \operatorname{tr}(\delta A\wedge F_A).
\]
Thus the Euler--Lagrange equation is
\[
  F_A=0 .
\]
On a manifold with boundary, the integration by parts also gives
\[
  \frac{k}{4\pi}\int_{\partial M}\operatorname{tr}(A\wedge\delta A),
\]
so a boundary condition, boundary action, or boundary phase space must be part
of the definition.

\section{Level Quantization}

For simply connected compact simple \(G\), choose the bilinear form so that
\[
  \frac{1}{24\pi^2}\int_{S^3}
  \operatorname{tr}\bigl(g^{-1}dg\bigr)^3\in\mathbb Z
\]
for maps \(g:S^3\to G\) representing the generator of \(\pi_3(G)\).  Under a
gauge transformation \(g:M\to G\), the Chern--Simons functional changes by
\[
  S_{\rm CS}(A^g)-S_{\rm CS}(A)
  =
  2\pi k\, n(g)
\]
on closed \(M\), where \(n(g)\in\mathbb Z\) is the winding number with the
chosen normalization.

\begin{proposition}[Gauge invariance of the exponentiated action]
\label{prop:cs-level-quantization}
With the integral normalization above, \(\exp(\ii S_{\rm CS})\) is invariant
under all gauge transformations on closed \(M\) if \(k\in\mathbb Z\).
\end{proposition}

\begin{proof}
The transformation formula gives
\(\exp(\ii S_{\rm CS}(A^g))=\exp(\ii S_{\rm CS}(A))\exp(2\pi\ii k n(g))\).
This equals \(\exp(\ii S_{\rm CS}(A))\) for all \(n(g)\in\mathbb Z\) exactly
when \(k\) is integral in the normalization used for the bilinear form.
\end{proof}

\section{Phase Space on a Surface}

Let \(M=\mathbb R\times\Sigma\).  Write \(A=A_0dt+A_\Sigma\).  The term
linear in time derivatives is
\[
  \frac{k}{4\pi}\int_{\mathbb R}dt
  \int_\Sigma \operatorname{tr}(A_\Sigma\wedge \partial_t A_\Sigma).
\]
The field \(A_0\) is a Lagrange multiplier imposing \(F_{A_\Sigma}=0\).
Therefore the classical phase space is the moduli space
\[
  \mathcal M_{\rm flat}(\Sigma,G)
  =
  \{A_\Sigma:F_{A_\Sigma}=0\}/\mathcal G_\Sigma
\]
with symplectic form
\[
  \omega_{\rm CS}
  =
  \frac{k}{4\pi}\int_\Sigma
  \operatorname{tr}(\delta A_\Sigma\wedge \delta A_\Sigma).
\]
Quantization of Chern--Simons theory assigns a finite-dimensional state space
to \(\Sigma\) in compact semisimple examples after the level and boundary
data are fixed.

\section{Wilson Lines and Framing}

For an oriented knot or link component \(K\subset M\) and a representation
\(R\) of \(G\), the Wilson line is
\[
  W_R(K)=\operatorname{Tr}_R\operatorname{Pexp}\int_K A .
\]
The quantum expectation value of Wilson links requires a framing: a
nonvanishing normal vector field along each component, or equivalently a
choice of point-splitting direction.  The framing dependence is a local
counterterm effect for the self-linking singularity of the line operator.

\section{Boundary Theory}

If \(M\) has boundary, gauge variation of the Chern--Simons action produces a
boundary term.  One way to restore a gauge-invariant variational problem is
to add boundary degrees of freedom whose action is the chiral Wess--Zumino--
Witten functional at level \(k\).  The boundary current algebra is the affine
Lie algebra determined by \(G\) and the same invariant form.  This is a
boundary reconstruction statement: the precise Hilbert space depends on the
chosen polarization and boundary condition.

\begin{openproblem}[Chern--Simons path integral as a QFT object]
\label{op:chern-simons-qft-object}
Give a regulator-level construction of Chern--Simons theory that simultaneously
controls gauge fixing, framing, Wilson-line renormalization, boundary
conditions, and functorial gluing.  Perturbative configuration-space
integrals, geometric quantization, and modular tensor category constructions
each capture part of this structure; a single local-QFT construction should
state the comparison maps among them.
\end{openproblem}
